% Options for packages loaded elsewhere
\PassOptionsToPackage{unicode}{hyperref}
\PassOptionsToPackage{hyphens}{url}
\PassOptionsToPackage{dvipsnames,svgnames,x11names}{xcolor}
\documentclass[
  10pt,
  fontset=fandol]{ctexart}
\usepackage{xcolor}
\usepackage[margin=16mm]{geometry}
\usepackage{amsmath,amssymb}
\setcounter{secnumdepth}{-\maxdimen} % remove section numbering
\usepackage{iftex}
\ifPDFTeX
  \usepackage[T1]{fontenc}
  \usepackage[utf8]{inputenc}
  \usepackage{textcomp} % provide euro and other symbols
\else % if luatex or xetex
  \usepackage{unicode-math} % this also loads fontspec
  \defaultfontfeatures{Scale=MatchLowercase}
  \defaultfontfeatures[\rmfamily]{Ligatures=TeX,Scale=1}
\fi
\usepackage{lmodern}
\ifPDFTeX\else
  % xetex/luatex font selection
\fi
% Use upquote if available, for straight quotes in verbatim environments
\IfFileExists{upquote.sty}{\usepackage{upquote}}{}
\IfFileExists{microtype.sty}{% use microtype if available
  \usepackage[]{microtype}
  \UseMicrotypeSet[protrusion]{basicmath} % disable protrusion for tt fonts
}{}
\makeatletter
\@ifundefined{KOMAClassName}{% if non-KOMA class
  \IfFileExists{parskip.sty}{%
    \usepackage{parskip}
  }{% else
    \setlength{\parindent}{0pt}
    \setlength{\parskip}{6pt plus 2pt minus 1pt}}
}{% if KOMA class
  \KOMAoptions{parskip=half}}
\makeatother
\setlength{\emergencystretch}{3em} % prevent overfull lines
\providecommand{\tightlist}{%
  \setlength{\itemsep}{0pt}\setlength{\parskip}{0pt}}
\usepackage{amsmath,amssymb,mathtools,needspace}
\xeCJKDeclareCharClass{CJK}{"2032,"0400->"04FF,"2160->"216B,"2460->"2473,"25A0,"25C7,"25CB,"25CF}
\IfFontExistsTF{Arial Unicode MS}{\xeCJKsetup{AutoFallBack=true}\setCJKfallbackfamilyfont{\CJKrmdefault}{Arial Unicode MS}}{}
\setcounter{secnumdepth}{0}
\setlength{\emergencystretch}{2em}
\allowdisplaybreaks
\usepackage{bookmark}
\IfFileExists{xurl.sty}{\usepackage{xurl}}{} % add URL line breaks if available
\urlstyle{same}
\hypersetup{
  pdftitle={向量和矩阵的范数},
  colorlinks=true,
  linkcolor={Maroon},
  filecolor={Maroon},
  citecolor={Blue},
  urlcolor={Blue},
  pdfcreator={LaTeX via pandoc}}

\title{向量和矩阵的范数}
\author{}
\date{}

\begin{document}
\maketitle

\subsection{7.5 向量和矩阵的范数}\label{ux5411ux91cfux548cux77e9ux9635ux7684ux8303ux6570}

为了研究线性方程组近似解的误差估计和迭代法的收敛性，需要对 \(\mathbf R^n\)（\(n\) 维向量空间）中向量（或 \(\mathbf R^{n\times n}\) 中矩阵）的``大小''引进某种度量------向量（或矩阵）范数的概念。向量范数概念是三维 Euclid 空间中向量长度概念的推广，在数值分析中起着重要作用。

我们用 \(\mathbf R^n\) 表示 \(n\) 维实向量空间，用 \(\mathbf C^n\) 表示 \(n\) 维复向量空间，首先将向量长度概念推广到 \(\mathbf R^n\)（或 \(\mathbf C^n\)）中。

\textbf{定义 7.1} 设

\[
x=(x_1,x_2,\cdots,x_n)^{\mathrm T},\quad y=(y_1,y_2,\cdots,y_n)^{\mathrm T}\in\mathbf R^n\text{（或 }\mathbf C^n\text{）},
\]

将实数 \((x,y)=y^{\mathrm T}x=\sum_{i=1}^{n}x_iy_i\)（或复数 \((x,y)=y^Hx=\sum_{i=1}^{n}x_i\overline{y_i}\)，其中 \(y^H=\overline{y^{\mathrm T}}\)）称为向量 \(x,y\) 的\textbf{数量积}。将非负实数

\[
\|x\|_2=(x,x)^{\frac12}=\left(\sum_{i=1}^{n}x_i^2\right)^{\frac12}
\]

或

\[
\|x\|_2=(x,x)^{\frac12}=\left(\sum_{i=1}^{n}|x_i|^2\right)^{\frac12}
\]

（定义续下页。）

\href{../../source-images/pdf-200.jpeg}{原页图像}

（续 7.5 定义 7.1。）

称为向量 \(x\) 的 \textbf{Euclid 范数}。

下述定理可在线性代数书中找到。

\textbf{定理 7.10} 设 \(x,y\in\mathbf R^n\)（或 \(\mathbf C^n\)），则

\(1^\circ\) \((x,x)=0\)，当且仅当 \(x=0\) 时成立；

\(2^\circ\) \((\alpha x,y)=\alpha(x,y)\)，\(\alpha\) 为实数（或 \((x,\alpha y)=\overline\alpha(x,y)\)，\(\alpha\) 为复数）；

\(3^\circ\) \((x,y)=(y,x)\)（或 \((x,y)=\overline{(y,x)}\)）；

\(4^\circ\) \((x_1+x_2,y)=(x_1,y)+(x_2,y)\)；

\(5^\circ\)（Cauchy-Schwarz 不等式）\(|(x,y)|\leqslant\|x\|_2\cdot\|y\|_2\)，等式当且仅当 \(x\) 与 \(y\) 线性相关时成立；

\(6^\circ\)（三角不等式）\(\|x+y\|_2\leqslant\|x\|_2+\|y\|_2\)。

还可以用其他办法来度量 \(\mathbf R^n\) 中向量的``大小''。例如 \(x=(x_1,x_2)^{\mathrm T}\in\mathbf R^2\)，用一个 \(x\) 的函数 \(N(x)=\max_{i=1,2}|x_i|\) 求度量 \(x\) 的``大小''，而且这种度量 \(x\)``大小''的方法计算起来比 Euclid 范数方便。在许多应用中，对度量向量 \(x\)``大小''的函数 \(N(x)\) 都要求是正定的、齐次的且满足三角不等式。下面给出向量范数的一般定义。

\textbf{定义 7.2（向量的范数）} 如果向量 \(x\in\mathbf R^n\)（或 \(\mathbf C^n\)）的某个实值函数 \(N(x)=\|x\|\)，满足条件

\[
\begin{aligned}
1^\circ\quad &\|x\|\geqslant0\quad(\|x\|=0\text{ 当且仅当 }x=0)\quad\text{（正定条件）};\\
2^\circ\quad &\|\alpha x\|=|\alpha|\|x\|,\quad\forall\alpha\in\mathbf R\text{（或 }\alpha\in\mathbf C\text{）};\\
3^\circ\quad &\|x+y\|\leqslant\|x\|+\|y\|\quad\text{（三角不等式）},
\end{aligned}
\tag{7.5.1}
\]

则称 \(N(x)\) 是 \(\mathbf R^n\)（或 \(\mathbf C^n\)）上的一个\textbf{向量范数}（或\textbf{模}）。

由条件 \(3^\circ\) 可推出不等式

\[
\big|\|x\|-\|y\|\big|\leqslant\|x-y\|.
\tag{7.5.2}
\]

下面给出几种常用的向量范数。

（1）向量的 \textbf{\(\infty\)-范数（最大范数）}：\(\|x\|_\infty=\max_{1\leqslant i\leqslant n}|x_i|\)。

容易验证这样定义的向量 \(x\) 的函数 \(N(x)=\|x\|_\infty\) 满足向量范数的三个条件。

（2）向量的 \textbf{1-范数}：\(\|x\|_1=\sum_{i=1}^{n}|x_i|\)。同样可证 \(N(x)=\|x\|_1\) 是 \(\mathbf R^n\) 上的一个向量范数。

（3）向量的 \textbf{2-范数}：\(\|x\|_2=(x,x)^{\frac12}=\left(\sum_{i=1}^{n}x_i^2\right)^{\frac12}\)。由定理 7.10 知，\(N(x)=\|x\|_2\) 是 \(\mathbf R^n\) 上一个向量范数，称为向量 \(x\) 的 Euclid 范数。

（4）向量的 \textbf{\(p\)-范数}：\(\|x\|_p=\left(\sum_{i=1}^{n}|x_i|^p\right)^{1/p}\)，其中 \(p\in[1,\infty)\)。可以证明向量函数 \(N(x)=\|x\|_p\) 是 \(\mathbf R^n\) 上向量的范数，且容易说明上述三种范数是 \(p\) 范数的特殊情况（\(\|x\|_\infty=\lim_{p\to\infty}\|x\|_p\)）。

\textbf{例 7.6} 计算向量 \(x=(1,-2,3)^{\mathrm T}\) 的各种范数。

（例 7.6 解续下页。）

\href{../../source-images/pdf-201.jpeg}{原页图像}

（续 7.5 例 7.6。）

\textbf{解} \(\|x\|_1=6,\|x\|_\infty=3,\|x\|_2=\sqrt{14}\)。

\textbf{定义 7.3} 设 \(\{x^{(k)}\}\) 为 \(\mathbf R^n\) 中一向量序列，\(x^*\in\mathbf R^n\)，记 \(x^{(k)}=(x_1^{(k)},x_2^{(k)},\cdots,x_n^{(k)})^{\mathrm T}\)，\(x^*=(x_1^*,x_2^*,\cdots,x_n^*)^{\mathrm T}\)。如果 \(\lim_{k\to\infty}x_i^{(k)}=x_i^*\)（\(i=1,2,\cdots,n\)），则称 \(x^{(k)}\) \textbf{收敛}于向量 \(x^*\)，记作

\[
\lim_{k\to\infty}x^{(k)}=x^*.
\]

\textbf{定理 7.11（\(N(x)\) 的连续性）} 设非负函数 \(N(x)=\|x\|\) 为 \(\mathbf R^n\) 上任一向量范数，则 \(N(x)\) 是 \(x\) 分量 \(x_1,x_2,\cdots,x_n\) 的连续函数。

\textbf{证明} 设 \(x=\sum_{i=1}^{n}x_ie_i,y=\sum_{i=1}^{n}y_ie_i\)，其中 \(e_i=(0,\allowbreak \cdots,\allowbreak 1,\allowbreak 0,\allowbreak \cdots,\allowbreak 0)^{\mathrm T}\)（即第 \(i\) 个元素为 1）。

只需证明当 \(x\to y\) 时 \(N(x)\to N(y)\) 即可。事实上，

\[
\begin{aligned}
|N(x)-N(y)|
&=\big|\|x\|-\|y\|\big|\leqslant\|x-y\|
=\left\|\sum_{i=1}^{n}(x_i-y_i)e_i\right\|\\
&\leqslant\sum_{i=1}^{n}|x_i-y_i|\|e_i\|
\leqslant\|x-y\|_\infty\sum_{i=1}^{n}\|e_i\|,
\end{aligned}
\]

即

\[
|N(x)-N(y)|\leqslant c\|x-y\|_\infty\to0\quad\text{（当 }x\to y\text{ 时）},
\]

其中

\[
c=\sum_{i=1}^{n}\|e_i\|.
\]

\textbf{定理 7.12（向量范数的等价性）} 设 \(\|x\|_s,\|x\|_t\) 为 \(\mathbf R^n\) 上向量的任意两种范数，则存在常数 \(c_1,c_2>0\)，使得

\[
c_1\|x\|_s\leqslant\|x\|_t\leqslant c_2\|x\|_s,\quad\text{对一切 }x\in\mathbf R^n.
\]

\textbf{证明} 只要就 \(\|x\|_s=\|x\|_\infty\) 证明上式成立即可，即证明存在常数 \(c_1,c_2>0\)，使

\[
c_1\leqslant\frac{\|x\|_t}{\|x\|_\infty}\leqslant c_2,\quad\text{对一切 }x\in\mathbf R^n\text{ 且 }x\neq0.
\]

考虑泛函 \(f(x)=\|x\|_t\geqslant0,\quad x\in\mathbf R^n\)。

记 \(S=\{x\mid\|x\|_\infty=1,x\in\mathbf R^n\}\)，则 \(S\) 是一个有界闭集。由于 \(f(x)\) 为 \(S\) 上的连续函数，所以 \(f(x)\) 于 \(S\) 上达到最大、最小值。设 \(x\in\mathbf R^n\) 且 \(x\neq0\)，则 \(\frac{x}{\|x\|_\infty}\in S\)，从而有

\[
f(x')=c_1\leqslant f\left(\frac{x}{\|x\|_\infty}\right)\leqslant c_2=f(x''),
\tag{7.5.3}
\]

其中 \(x',x''\in S\)。显然 \(c_1,c_2>0\)，上式为 \(c_1\leqslant\left\|\frac{x}{\|x\|_\infty}\right\|\leqslant c_2\)，即

\[
c_1\|x\|_\infty\leqslant\|x\|_t\leqslant c_2\|x\|_\infty,\quad\text{对一切 }x\in\mathbf R^n.
\]

注意，定理 7.12 不能推广到无穷维空间。由定理 7.12 可得到结论：如果在一种范数意义下向量序列收敛，则在任何一种范数意义下该向量序列亦收敛。

\href{../../source-images/pdf-202.jpeg}{原页图像}

（续 7.5 向量和矩阵的范数。）

\textbf{定理 7.13} \(\lim_{k\to\infty}x^{(k)}=x^*\Longleftrightarrow\|x^{(k)}-x^*\|\to0\)（当 \(k\to\infty\) 时），其中 \(\|\cdot\|\) 为向量的任一种范数。

\textbf{证明} 显然，\(\lim_{k\to\infty}x^{(k)}=x^*\Longleftrightarrow\|x^{(k)}-x^*\|_\infty\to0\)（当 \(k\to\infty\) 时），而对于 \(\mathbf R^n\) 上任一种范数 \(\|\cdot\|\)，由定理 7.11，存在常数 \(c_1,c_2>0\)，使

\[
c_1\|x^{(k)}-x^*\|_\infty\leqslant\|x^{(k)}-x^*\|\leqslant c_2\|x^{(k)}-x^*\|_\infty,
\]

于是又有

\[
\|x^{(k)}-x^*\|_\infty\to0\Longleftrightarrow\|x^{(k)}-x^*\|\to0\quad\text{（当 }k\to\infty\text{ 时）}.
\]

下面将向量范数概念推广到矩阵上去。用 \(\mathbf R^{n\times n}\) 表示 \(n\times n\) 矩阵的集合（本质上是和 \(\mathbf R^{n\times n}\) 一样的向量空间），则由 \(\mathbf R^{n\times n}\) 上 2-范数可以得到 \(\mathbf R^{n\times n}\) 中矩阵的一种范数

\[
F(A)=\|A\|_F=\left(\sum_{i,j=1}^{n}a_{ij}^{2}\right)^{\frac12}
\]

称为 \(A\) 的 \textbf{Frobenius 范数}。\(\|A\|_F\) 显然满足正定性、齐次性及三角不等式。

下面给出矩阵范数的一般定义。

\textbf{定义 7.4（矩阵范数）} 如果矩阵 \(A\in\mathbf R^{n\times n}\) 的某个非负的实值函数 \(N(A)=\|A\|\)，满足条件

\[
\begin{aligned}
1^\circ\quad &\|A\|\geqslant0\quad(\|A\|=0\Longleftrightarrow A=0)\quad\text{（正定条件）};\\
2^\circ\quad &\|cA\|=|c|\|A\|,\quad c\text{ 为实数}\quad\text{（齐次条件）};\\
3^\circ\quad &\|A+B\|\leqslant\|A\|+\|B\|\quad\text{（三角不等式）};\\
4^\circ\quad &\|AB\|\leqslant\|A\|\|B\|,
\end{aligned}
\tag{7.5.4}
\]

则称 \(N(A)\) 是 \(\mathbf R^{n\times n}\) 上的一个\textbf{矩阵范数}（或\textbf{模}）。

上面定义的 \(F(A)=\|A\|_F\) 就是 \(\mathbf R^{n\times n}\) 上的一个矩阵范数。

由于在大多数与估计有关的问题中，矩阵和向量会同时参与讨论，所以希望引进一种矩阵的范数，它是和向量范数相联系而且和向量范数相容的，即

\[
\|Ax\|\leqslant\|A\|\|x\|
\tag{7.5.5}
\]

对于任意向量 \(x\in\mathbf R^n\) 及 \(A\in\mathbf R^{n\times n}\) 都成立。

为此再引进一种矩阵范数。

\textbf{定义 7.5（矩阵的算子范数）} 设 \(x\in\mathbf R^n,A\in\mathbf R^{n\times n}\)，给出一种向量范数 \(\|x\|_v\)（如 \(v=1,2\) 或 \(\infty\)），相应地定义一个矩阵的非负函数

\[
\|A\|_v=\max_{x\neq0}\frac{\|Ax\|_v}{\|x\|_v}.
\tag{7.5.6}
\]

可验证 \(\|A\|_v\) 满足定义 7.4（见下面定理），所以 \(\|A\|_v\) 是 \(\mathbf R^{n\times n}\) 上的一个矩阵范数，称为 \(A\) 的\textbf{算子范数}。

\textbf{定理 7.14} 设 \(\|x\|_v\) 是 \(\mathbf R^n\) 上一个向量范数，则 \(\|A\|_v\) 是 \(\mathbf R^{n\times n}\) 上的矩阵范数，且满足相容条件

\[
\|Ax\|_v\leqslant\|A\|_v\|x\|_v.
\tag{7.5.7}
\]

\textbf{证明} 由式（7.5.6）知，相容性条件式（7.5.7）是显然的。现只验证定义 7.4 中条

（证明接下页。）

\begin{quote}
校注（agent 补充，原书疑误）：定理 7.13 证明中原文引用``定理 7.11''，这里保留；所用的两侧常数比较直接对应上一页的定理 7.12（向量范数的等价性），定理 7.11 是连续性定理。\href{../assets/ch07b/202-reference.png}{原文局部图}。
\end{quote}

\href{../../source-images/pdf-203.jpeg}{原页图像}

（续 7.5 定理 7.14 证明。）

件 \(4^\circ\)。

由式（7.5.7），有

\[
\|ABx\|_v\leqslant\|A\|_v\|Bx\|_v\leqslant\|A\|_v\|B\|_v\|x\|_v.
\]

当 \(x\neq0\) 时，有

\[
\frac{\|ABx\|_v}{\|x\|_v}\leqslant\|A\|_v\|B\|_v,\quad
\|AB\|_v=\max_{x\neq0}\frac{\|ABx\|_v}{\|x\|_v}\leqslant\|A\|_v\|B\|_v.
\]

显然这种矩阵范数 \(\|A\|_v\) 依赖于向量范数 \(\|x\|_v\) 的具体含义。也就是说，当给出一种具体的向量范数 \(\|x\|_v\) 时，相应地就得到了一种矩阵范数 \(\|A\|_v\)。

\textbf{定理 7.15} 设 \(x\in\mathbf R^n,A\in\mathbf R^{n\times n}\)，则

\(1^\circ\) \(\|A\|_\infty=\max_{1\leqslant i\leqslant n}\sum_{j=1}^{n}|a_{ij}|\)（称为 \(A\) 的\textbf{行范数}）；

\(2^\circ\) \(\|A\|_1=\max_{1\leqslant j\leqslant n}\sum_{i=1}^{n}|a_{ij}|\)（称为 \(A\) 的\textbf{列范数}）；

\(3^\circ\) \(\|A\|_2=\sqrt{\lambda_{\max}(A^{\mathrm T}A)}\)（称为 \(A\) 的 \textbf{2-范数}），

其中 \(\lambda_{\max}(A^{\mathrm T}A)\) 表示 \(A^{\mathrm T}A\) 的最大特征值。

\textbf{证明} 只就 \(1^\circ\)、\(3^\circ\) 给出证明，\(2^\circ\) 同理可证。

\(1^\circ\) 设 \(x=(x_1,x_2,\cdots,x_n)^{\mathrm T}\neq0\)，不妨设 \(A\neq0\)。记 \(t=\max_i|x_i|\)，\(\mu=\max_{1\leqslant i\leqslant n}\sum_{j=1}^{n}|a_{ij}|\)，则

\[
\|Ax\|_\infty
=\max_{1\leqslant i\leqslant n}\left|\sum_{j=1}^{n}a_{ij}x_j\right|
\leqslant\max_i\sum_{j=1}^{n}|a_{ij}||x_j|
\leqslant t\max_i\sum_{j=1}^{n}|a_{ij}|.
\]

这说明，对于任何非零 \(x\in\mathbf R^n\)，有

\[
\frac{\|Ax\|_\infty}{\|x\|_\infty}\leqslant\mu.
\tag{7.5.8}
\]

下面说明，有一向量 \(x_0\neq0\)，使 \(\frac{\|Ax_0\|_\infty}{\|x_0\|_\infty}=\mu\)。设 \(\mu=\sum_{j=1}^{n}|a_{i_0j}|\)，取向量

\[
x_0=(x_1,x_2,\cdots,x_n)^{\mathrm T},
\]

其中

\[
x_j=\operatorname{sgn}(a_{i_0j})\quad(j=1,2,\cdots,n).
\]

显然 \(\|x_0\|_\infty=1\)，且 \(Ax_0\) 的第 \(i_0\) 个分量为 \(\sum_{i=1}^{n}a_{i_0j}x_j=\sum_{j=1}^{n}|a_{i_0j}|\)，这说明

\[
\|Ax_0\|_\infty=\max_i\sum_{j=1}^{n}|a_{ij}x_j|=\sum_{j=1}^{n}|a_{i_0j}|=\mu.
\]

\(3^\circ\) 由于 \(\|Ax\|_2^2=(Ax,Ax)=(A^{\mathrm T}Ax,x)\geqslant0\)，对于一切 \(x\in\mathbf R^n\)，从而 \(A^{\mathrm T}A\) 的特征值为非负实数，设为

\[
\lambda_1\geqslant\lambda_2\geqslant\cdots\geqslant\lambda_n\geqslant0.
\tag{7.5.9}
\]

\(A^{\mathrm T}A\) 为对称矩阵，设 \(u_1,u_2,\cdots,u_n\) 为 \(A^{\mathrm T}A\) 的相应于式（7.5.9）的特征向量且

（证明续下页。）

\begin{quote}
校注（agent 补充，原书疑误）：\(Ax_0\) 第 \(i_0\) 个分量表达式的第一个求和号下限原图为 \(i=1\)，这里保留；该和式的项是 \(a_{i_0j}x_j\)，求和指标应为 \(j\)，相应下限应为 \(j=1\)。\href{../assets/ch07b/203-sum-index.png}{原式局部图}。
\end{quote}

\href{../../source-images/pdf-204.jpeg}{原页图像}

（续 7.5 定理 7.15 证明。）

\((u_i,u_j)=\delta_{ij}\)，又设 \(x\in\mathbf R^n\) 为任一非零向量，于是有

\[
x=\sum_{i=1}^{n}c_iu_i,
\]

其中 \(c_i\) 为组合系数。

\[
\frac{\|Ax\|_2^2}{\|x\|_2^2}
=\frac{(A^{\mathrm T}Ax,x)}{(x,x)}
\leqslant\frac{\sum_{i=1}^{n}c_i^2\lambda_1}{\sum_{i=1}^{n}c_i^2}=\lambda_1.
\]

另一方面，取 \(x=u_1\)，则上式等号成立，故

\[
\|A\|_2=\max_{x\neq0}\frac{\|Ax\|_2}{\|x\|_2}=\sqrt{\lambda_1}=\sqrt{\lambda_{\max}(A^{\mathrm T}A)}.
\]

由定理 7.15 看出，计算一个矩阵的 \(\|A\|_\infty,\|A\|_1\) 还是比较容易的，而矩阵的 2-范数 \(\|A\|_2\) 在计算上不方便。但是矩阵的 2-范数具有许多好的性质，它在理论上是有用的。

\textbf{例 7.7} \(A=\begin{pmatrix}1&-2\\-3&4\end{pmatrix}\)，计算 \(A\) 的各种范数。

\textbf{解} \(\|A\|_1=6,\|A\|_\infty=7,\|A\|_F\approx5.477,\|A\|_2=\sqrt{15+\sqrt{221}}\approx5.46\)。

对于复矩阵（即 \(A\in\mathbf C^{n\times n}\)），定理 7.15 中的 \(1^\circ\)、\(2^\circ\) 显然成立；\(3^\circ\) 则应改为

\[
\|A\|_2=\max_{x\neq0}\left(\frac{x^HA^HAx}{x^Hx}\right)=\sqrt{\lambda_{\max}(A^HA)}.
\]

其中 \(A^H=\overline A^{\mathrm T}\)（\(A\) 共轭转置）。

\textbf{定义 7.6} 设 \(A\in\mathbf R^{n\times n}\) 的特征值为 \(\lambda_i\)（\(i=1,2,\cdots,n\)），称 \(\rho(A)=\max_{1\leqslant i\leqslant n}|\lambda_i|\) 为 \(A\) 的\textbf{谱半径}。

\textbf{定理 7.16（特征值上界）} 设 \(A\in\mathbf R^{n\times n}\)，则 \(\rho(A)\leqslant\|A\|\)，即 \(A\) 的谱半径不超过 \(A\) 的任何一种算子范数（对于 \(\|A\|_F\) 亦对）。

\textbf{证明} 设 \(\lambda\) 是 \(A\) 的任一特征值，\(x\) 为相应的特征向量，则 \(Ax=\lambda x\)，由式（7.5.7）得

\[
|\lambda|\|x\|=\|\lambda x\|=\|Ax\|\leqslant\|A\|\|x\|,
\]

即

\[
|\lambda|\leqslant\|A\|.
\]

\textbf{定理 7.17} 如果 \(A\in\mathbf R^{n\times n}\) 为对称矩阵，则 \(\|A\|_2=\rho(A)\)。

证明留作习题。

\textbf{定理 7.18} 如果 \(\|B\|<1\)，则 \(I\pm B\) 为非奇异矩阵，且 \(\|(I\pm B)^{-1}\|\leqslant\frac1{1-\|B\|}\)，其中 \(\|\cdot\|\) 是指矩阵的算子范数。

\textbf{证明} 用反证法。若 \(\det(I\pm B)=0\)，则 \((I\pm B)x=0\) 有非零解，即存在 \(x_0\neq0\) 使

（证明续下页。）

\begin{quote}
校注（agent 补充，原书疑误）：复矩阵 2-范数公式的 Rayleigh 商外，原图没有平方根或 \(1/2\) 次幂，此处照录。该最大商等于 \(\lambda_{\max}(A^HA)=\|A\|_2^2\)；要与两侧 \(\|A\|_2\) 和 \(\sqrt{\lambda_{\max}(A^HA)}\) 一致，中间项需开平方。\href{../assets/ch07b/204-complex-norm.png}{原式局部图}。
\end{quote}

\href{../../source-images/pdf-205.jpeg}{原页图像}

（续 7.5 定理 7.18 证明。）

\(Bx_0=x_0\)，\(\frac{\|Bx_0\|}{\|x_0\|}=1\)，故 \(\|B\|\geqslant1\)，与假设矛盾。又由 \((I\pm B)(I\pm B)^{-1}=I\)，有

\[
(I\pm B)^{-1}=I\mp B(I-B)^{-1},
\]

从而

\[
\|(I\pm B)^{-1}\|\leqslant\|I\|+\|B\|\|(I\pm B)^{-1}\|,
\]

\[
\|(I\pm B)^{-1}\|\leqslant\frac1{1-\|B\|}.
\]

\begin{center}\rule{0.5\linewidth}{0.5pt}\end{center}

\subsection{相关原页校注（agent 补充）}\label{ux76f8ux5173ux539fux9875ux6821ux6ce8agent-ux8865ux5145}

以下校注来自本节覆盖的原页，可能也涉及同页相邻小节；原文照录与校正说明分开保留。

\subsubsection{原书 PDF 页 205 的校注}\label{ux539fux4e66-pdf-ux9875-205-ux7684ux6821ux6ce8}

\href{../pages/pdf-205.md}{完整原页转录} · \href{../../source-images/pdf-205.jpeg}{原页图像}

校注记录：ch07b-E007。

\begin{quote}
校注（agent 补充，原书疑误）：本页定理 7.18 证明原文为 \(Bx_0=x_0\)，以及 \((I\pm B)^{-1}=I\mp B(I-B)^{-1}\)，均按印刷内容保留。若同时处理 \(I+B\) 和 \(I-B\) 两种情形，前式应为 \(Bx_0=\mp x_0\)，后式右端的逆矩阵应为 \((I\pm B)^{-1}\)；原文写法只与减号分支一致。前式符号不影响随后所取范数等于 1 的结论。\href{../assets/ch07b/205-inverse-signs.png}{原式局部图}。
\end{quote}

\subsection{本节覆盖原页的校注记录}\label{ux672cux8282ux8986ux76d6ux539fux9875ux7684ux6821ux6ce8ux8bb0ux5f55}

下列记录按原页关联，含同页相邻小节的校注；计算时同时读取上文校注和完整原页。

\begin{itemize}
\tightlist
\item
  \texttt{ch07b-E004}：原书 PDF 页 202
\item
  \texttt{ch07b-E005}：原书 PDF 页 203
\item
  \texttt{ch07b-E006}：原书 PDF 页 204
\item
  \texttt{ch07b-E007}：原书 PDF 页 205
\end{itemize}

\end{document}
